\chapter{Gestione delle sessioni}

Alla nostra applicazione manca ancora una funzionalità essenziale: più utenti
devono poter registrarsi, autenticarsi e gestire ciascuno la \textbf{propria}
lista. Realizzarla richiede un concetto nuovo, il \term{tracciamento della
sessione}.

\section{Il problema}

Torniamo al punto lasciato in sospeso nel primo capitolo: HTTP è
\textbf{stateless}. Una richiesta, di per sé, non contiene alcuna informazione
sulle richieste precedenti.

Come fa il server a sapere che un utente ha già effettuato il login?

\begin{figure}[H]
  \centering
  \begin{tikzpicture}[font=\Lato\small,
      hdr/.style={wtnode, minimum width=3.0cm, minimum height=0.9cm},
      life/.style={draw=neutral!45!white, line width=0.7pt, dashed},
      msg/.style={-{Stealth[length=6pt,width=5pt]}, line width=0.9pt,
                  draw=primary!70!black},
      rsp/.style={-{Stealth[length=6pt,width=5pt]}, line width=0.9pt,
                  draw=accent!70!black},
      ml/.style={font=\FiraCodeNL\scriptsize, midway, above, inner sep=2.5pt,
                 fill=white, text=neutral!15!black}]

    \def\xb{0} \def\xs{9.0}

    \node[hdr] (B) at (\xb,0) {\textbf{Utente}};
    \node[hdr] (S) at (\xs,0) {\textbf{Application Server}};
    \draw[life] (B.south) -- (\xb,-4.0);
    \draw[life] (S.south) -- (\xs,-4.0);

    \draw[msg] (\xb,-0.9) -- (\xs,-0.9) node[ml]{POST /login usr=luigi\&pwd=pass};
    \draw[rsp] (\xs,-1.8) -- (\xb,-1.8) node[ml]{200 OK};
    \draw[msg] (\xb,-2.7) -- (\xs,-2.7) node[ml]{GET /todo};

    \node[font=\Lato\Large\bfseries, text=red-gradient-6] at (\xs+0.9,-3.3) {?};
    \node[font=\Lato\scriptsize\itshape, text=red-gradient-6!80!black,
          anchor=west, align=left] at (\xs+1.3,-3.3)
      {chi è questo\\ utente?};
  \end{tikzpicture}
  \caption{Senza stato, la terza richiesta arriva da uno sconosciuto.}
  \label{fig:stateless-problem}
\end{figure}

\dfn{Tracciamento della sessione}{
  Il \term{tracciamento della sessione} (\textit{session tracking}) consiste nel
  mantenere informazioni \textbf{con stato} sulle interazioni di un utente
  attraverso più richieste HTTP. Permette ai server di \guillemotleft riconoscere\guillemotright\ gli
  utenti e le loro azioni, e quindi di adattare le risposte a utenti e
  situazioni diverse.}

\section{Un approccio ingenuo: i cookie}

Un modo che già conosciamo perché il server \guillemotleft memorizzi\guillemotright\ dati fra
richieste diverse è impostare dei \term{cookie}. L'idea è la seguente.

\begin{enumerate}
  \item L'utente invia una richiesta a una pagina dinamica, con nome utente e
        password come parametri.
  \item La pagina verifica se le credenziali sono corrette.
    \begin{itemize}
      \item Se lo sono, il server imposta dei cookie per tenere traccia
            dell'interazione, per esempio
            \codel{Set-Cookie: auth=true; username=luigi}.
      \item Se non lo sono, mostra una pagina di errore e non imposta alcun
            cookie.
    \end{itemize}
  \item Le pagine dinamiche che richiedono autenticazione controllano il cookie
        \code{auth}: se è impostato mostrano il contenuto, altrimenti
        reindirizzano alla pagina di login.
\end{enumerate}

\subsection{Gestire il login}

\begin{lstlisting}[style=js, name=login.js]
if (isAuthenticated) {
  response.writeHead(300, {
    "Location": "/",
    "Set-Cookie": [
      `auth=true; max-age=${60 * 60}`,
      `username=${user.username}; max-age=${60 * 60}`
    ]
  });
  response.end();
}
\end{lstlisting}

Si impostano due cookie --- uno con il nome utente, l'altro un semplice valore
booleano --- entrambi con scadenza a un'ora, e si reindirizza alla homepage.

Se invece le credenziali non sono corrette, si mostra di nuovo la pagina di
login restituendo un codice \code{401}:

\begin{lstlisting}[style=js, name=login.js]
if (!isAuthenticated) {
  let renderedContent = pug.renderFile("./templates/login.pug",
    { "error": "Autenticazione fallita. Controlla le credenziali." });
  response.writeHead(401, { "Content-Type": "text/html" });
  response.end(renderedContent);
}
\end{lstlisting}

\subsection{Recuperare le informazioni di sessione}

Vogliamo gestire le richieste in modo diverso a seconda che l'utente sia
autenticato o meno: se un utente non autorizzato prova a visitare \code{/todo}
gli serviamo una pagina di errore con codice \code{401}, mentre un utente
autenticato vede la propria lista.

\begin{lstlisting}[style=js, name=auth.js]
function checkUserAuthentication(request) {
  let cookies = parseCookies(request);
  if (cookies.auth) {
    return [true, cookies.username];
  } else {
    return [false, undefined];
  }
}
\end{lstlisting}

Purtroppo, anche per analizzare i cookie siamo da soli. L'header \code{Cookie}
si presenta così: \codel{Cookie: auth=true; username=luigi}.

\begin{lstlisting}[style=js, name=cookies.js]
function parseCookies(request) {
  const list = {};
  const cookieHeader = request.headers?.cookie;
  if (!cookieHeader) return list;

  cookieHeader.split(`;`).forEach(function (cookie) {
    let [name, ...rest] = cookie.split(`=`);
    name = name?.trim();
    if (!name) return;

    const value = rest.join(`=`).trim();
    if (!value) return;

    list[name] = decodeURIComponent(value);
  });

  return list;
}
\end{lstlisting}

\info{Perché quel \code{rest.join}}{
  Il codice cerca di essere robusto quando il \textbf{valore} di un cookie
  contiene a sua volta un \code{=}. Spezzare su \code{=} produce più di due
  parti, e ricomporre tutto ciò che segue il primo \code{=} è l'unico modo
  corretto. Di norma i cookie sono codificati come URL, ma la specifica non lo
  impone.}

\subsection{Rendere disponibile lo stato ai controller}

Tutte le pagine della nostra applicazione sono influenzate dall'autenticazione:
la barra di navigazione, per esempio, mostrerà \guillemotleft Benvenuto, utente\guillemotright\
invece di \guillemotleft Login\guillemotright. Vogliamo quindi che questa informazione sia
disponibile a tutti i controller.

Un modo è calcolarla una volta sola in \code{handleRequest} e passarla lungo
la catena come oggetto \term{context}.

\begin{lstlisting}[style=js, name=router.js]
function handleRequest(request, response) {
  let [isUserAuthenticated, username] = checkUserAuthentication(request);

  let context = {
    isUserAuthenticated: isUserAuthenticated,
    username: username
  };

  switch (request.url) {
    /* ... altre rotte ... */
    case "/todo":
      handleTodoListRequest(request, response, context); break;
  }
}
\end{lstlisting}

\begin{lstlisting}[style=js, name=todo-controller.js]
function handleTodoListRequest(request, response, context = {}) {
  if (!context.isUserAuthenticated) {
    handleError(request, response, 401, "Non autorizzato!");
  } else {
    switch (request.method) {
      case "GET":
        handleTodoListRequestGet(request, response, context); break;
      case "POST":
        handleTodoListRequestPost(request, response, context); break;
      default:
        handleError(request, response, 405, "Metodo non supportato");
    }
  }
}
\end{lstlisting}

\section{I problemi dell'approccio a cookie}

L'approccio funziona ed è semplice da implementare. Purtroppo soffre di un
problema \textbf{critico}.

\error{I cookie sono interamente controllati dal client}{
  Manipolarli è banale: bastano i DevTools. Un utente può cambiare il valore del
  proprio cookie \code{username} in qualunque altro nome utente, oppure
  impostarsi da solo un cookie \code{auth=true} \textbf{senza mai passare dalla
  pagina di login}. La nostra autenticazione diventa del tutto inutile.

  Il problema della manomissione si può mitigare con i \term{cookie firmati}: il
  server firma i cookie con la propria chiave privata e, finché la chiave resta
  segreta, riconosce quelli manomessi. Resta però il fatto che i dati viaggiano
  sul client.}

Ci sono poi limiti pratici. I browser impongono restrizioni sui cookie per
evitare problemi di prestazioni:

\begin{itemize}
  \item la dimensione di un cookie è generalmente limitata a \textbf{4096 byte};
  \item ogni host può memorizzare in genere dai \textbf{20 ai 50 cookie} per
        dominio.
\end{itemize}

E se dovessimo memorizzare più dati di sessione?

\section{Le sessioni web}

\dfn{Sessione web}{
  Quando serve tenere traccia di informazioni attraverso più richieste, il server
  crea una \term{sessione} per un dato client. Una sessione è una struttura dati
  che memorizza coppie chiave-valore, \textbf{gestita e conservata sul server}.
  È identificata da un \term{identificatore di sessione} univoco (o
  \textit{session token}), che viene passato al client, generalmente come
  cookie.}

Il punto chiave è che i \textbf{dati restano sul server}: il client possiede
soltanto un identificatore opaco, e non può metterci mano.

\begin{figure}[H]
  \centering
  \begin{tikzpicture}[font=\Lato\small,
      hdr/.style={wtnode, minimum width=3.0cm, minimum height=0.9cm},
      life/.style={draw=neutral!45!white, line width=0.7pt, dashed},
      msg/.style={-{Stealth[length=6pt,width=5pt]}, line width=0.9pt,
                  draw=primary!70!black},
      rsp/.style={-{Stealth[length=6pt,width=5pt]}, line width=0.9pt,
                  draw=accent!70!black},
      ml/.style={font=\FiraCodeNL\scriptsize, midway, above, inner sep=2.5pt,
                 fill=white, text=neutral!15!black}]

    \def\xb{0} \def\xs{8.4}

    \node[hdr] (B) at (\xb,0) {\textbf{Utente A}};
    \node[hdr] (S) at (\xs,0) {\textbf{Application Server}};
    \draw[life] (B.south) -- (\xb,-3.4);
    \draw[life] (S.south) -- (\xs,-3.4);

    \draw[msg] (\xb,-0.9) -- (\xs,-0.9) node[ml]{POST /login usr=john\&pwd=pass};
    \draw[rsp] (\xs,-1.7) -- (\xb,-1.7) node[ml]{200 OK Set-Cookie: SID=ahs37...435};
    \draw[msg] (\xb,-2.5) -- (\xs,-2.5) node[ml]{GET /todos Cookie: SID=ahs37...435};
    \draw[rsp] (\xs,-3.2) -- (\xb,-3.2) node[ml]{200 OK lista di john};

    % store lato server
    \node[wtnodeplain, anchor=north west, align=left,
          font=\FiraCodeNL\scriptsize] (st) at (10.4,0.4)
      {ahs37...435\\ \ \ auth: true\\ \ \ username: john\\[3pt]
       4fg23...43u\\ \ \ auth: true\\ \ \ username: billy};
    \node[font=\Lato\scriptsize, text=green-gradient-6!60!black, anchor=south west]
      at (10.4,0.5) {dati di sessione (sul server)};
    \begin{scope}[on background layer]
      \node[rounded corners=4pt, draw=green-gradient-6, line width=0.9pt,
            fill=green-gradient-1!35!white, fit=(st), inner sep=7pt] {};
    \end{scope}
  \end{tikzpicture}
  \caption{Il client possiede solo l'identificatore; i dati restano sul server.}
  \label{fig:sessions}
\end{figure}

\section{Implementare le sessioni da zero}

Realizzare un meccanismo di sessione di base in Node.js è piuttosto semplice.
Si possono usare gli oggetti \code{Map} di JavaScript: una prima mappa associa
ciascun identificatore di sessione a una sessione, che è a sua volta una mappa
di coppie chiave-valore.

\begin{lstlisting}[style=js, name=Session.js]
class Session {

  constructor() {
    this.sessionStore = new Map();
  }

  createSession() {
    let sessionId = this.generateNewSessionId();
    this.sessionStore.set(sessionId, new Map());
    return sessionId;
  }

  storeSessionData(sessionId, key, value) {
    return this.getSessionById(sessionId)?.set(key, value);
  }

  getSessionFromRequest(request) {
    let requestSessionId = parseCookies(request)["sessionId"];
    return this.getSessionById(requestSessionId);
  }

  /* altri metodi ... */
}
\end{lstlisting}

\subsection{Usare le sessioni}

Quando un utente effettua il login, creiamo una nuova sessione, vi memorizziamo
i dati rilevanti e passiamo l'identificatore come cookie.

\begin{lstlisting}[style=js, name=login.js]
if (isAuthenticated) {
  // crea una nuova sessione per l'utente
  let sessionId = session.createSession();
  session.storeSessionData(sessionId, "username", user.username);
  session.storeSessionData(sessionId, "auth", true);

  response.writeHead(300, {
    "Location": "/",
    "Set-Cookie": [
      `sessionId=${sessionId}; max-age=${60 * 60}`
    ]
  });
  response.end();
}
\end{lstlisting}

Quando serve accedere ai dati di sessione:

\begin{lstlisting}[style=js, name=router.js]
function handleRequest(request, response) {
  let userSession = session.getSessionFromRequest(request);

  let context = {
    isUserAuthenticated: userSession?.get("auth"),
    username: userSession?.get("username")
  };

  /* resto del codice */
}
\end{lstlisting}

Il codice del router è praticamente identico a quello della versione con i
cookie: cambia solo \textbf{da dove} arrivano i dati. Il cookie ora contiene un
identificatore opaco, e manometterlo non serve a nulla: se l'identificatore non
corrisponde a una sessione esistente sul server, \code{getSessionById}
restituisce \code{undefined}.

\warning{Le sessioni in memoria non bastano in produzione}{
  Conservare le sessioni in una \code{Map} in memoria funziona benissimo per
  imparare, ma ha due limiti concreti. Al \textbf{riavvio} del server tutte le
  sessioni spariscono e tutti gli utenti devono rifare il login. E se
  l'applicazione gira su \textbf{più macchine} dietro un bilanciatore di carico,
  una richiesta servita dalla macchina B non trova la sessione creata sulla
  macchina A. In produzione le sessioni si conservano in un archivio condiviso
  (Redis, un database), oppure si adottano token autoconsistenti e firmati.}

\subsection{Altri metodi di tracciamento}

Usare un cookie per memorizzare l'identificatore di sessione è l'approccio più
diffuso: è semplice e ampiamente supportato. Ne esistono però altri:

\begin{itemize}
  \item \term{URL rewriting}: l'identificatore viene aggiunto a ogni URL come
        parametro di query;
  \item \term{campi form nascosti}: l'identificatore è memorizzato in campi
        nascosti dei form e viene inviato insieme agli altri dati.
\end{itemize}

Entrambi hanno lo svantaggio di esporre l'identificatore in luoghi dove
è facile che finisca per essere condiviso o registrato nei log: si usano
soprattutto quando i cookie non sono disponibili.

\section{Riferimenti}

\begin{itemize}
  \item \textbf{Introduction to Web Applications Development}, Carles Mateu.
        Liberamente disponibile su \textit{archive.org} sotto licenza GNU FDL. \\
        Parte rilevante: sezione 3.11 (\textit{Session monitoring}); si possono
        ignorare gli esempi con le Java Servlet.
  \item \textbf{Using HTTP cookies} --- MDN web docs. \\
        \urlx{https://developer.mozilla.org/en-US/docs/Web/HTTP/Cookies}
  \item \textbf{RFC 6265: HTTP State Management Mechanism}. \\
        \urlx{https://datatracker.ietf.org/doc/html/rfc6265}
\end{itemize}
